% !TeX program = xelatex
\documentclass[11pt,a4paper]{article}
\usepackage[margin=23mm]{geometry}
\usepackage{fontspec}
\usepackage{polyglossia}
\setdefaultlanguage{russian}
\setotherlanguage{english}
\setmainfont{DejaVu Serif}
\setsansfont{DejaVu Sans}
\newfontfamily\cyrillicfont{DejaVu Serif}
\newfontfamily\cyrillicfontsf{DejaVu Sans}
\usepackage{amsmath,amssymb,mathtools,graphicx,xcolor,booktabs,enumitem,hyperref}
\definecolor{NESblue}{RGB}{5,67,114}
\definecolor{NESred}{RGB}{163,22,70}
\definecolor{NESgray}{RGB}{80,78,83}
\hypersetup{unicode=true,colorlinks=true,linkcolor=NESblue,urlcolor=NESblue,
 pdftitle={Микроэкономика 1: полный текст для преподавателя},pdfauthor={Арсений Столяров}}
\setlength{\parindent}{0pt}
\setlength{\parskip}{.55em}
\setlength{\emergencystretch}{3em}
\raggedright
\raggedbottom
\allowdisplaybreaks
\setcounter{tocdepth}{1}
\newcommand{\slide}[2]{\subsection*{\textcolor{NESblue}{Слайд #1.} #2}\pdfbookmark[2]{Слайд #1. #2}{script-#1}\hypertarget{script-#1}{}}
\newcommand{\cue}[1]{{\small\color{NESgray}\textit{Для преподавателя: #1}\par}}
\newcommand{\ask}[1]{\par\textcolor{NESred}{\textbf{Вопрос аудитории.}} #1\par}
\newcommand{\answer}[1]{\par\textbf{Разбор.} #1\par}
\newcommand{\reserve}{\cue{Дополнительный материал. Если времени мало, можно сократить или пропустить.}}
\newcommand{\pictureof}[1]{\begin{center}\includegraphics[width=.92\linewidth]{assets/script/slide-#1.pdf}\end{center}}
\newcommand{\R}{\mathbb R}
\begin{document}
\begin{titlepage}
\begin{flushright}\includegraphics[width=32mm]{assets/nes-english.pdf}\end{flushright}
\vspace{15mm}
{\Huge\sffamily\bfseries\color{NESblue}Микроэкономика\par}
\vspace{7mm}
{\Large Полный текст первой лекции\\для подготовки преподавателя\par}
\vspace{12mm}
{\large Арсений Столяров\par}
PhD in Economics, University of Wisconsin--Madison
\vspace{8mm}

Искусственный интеллект и финансовые технологии\\
Российская экономическая школа\\
27 октября 2026 года
\vfill
\textbf{Соответствует текущей презентации:} 54 основных слайда и 10 слайдов приложения. Номера ниже --- логические номера слайдов, совпадающие со страницами раздатки; в PDF с анимацией одному слайду соответствует несколько страниц.

Текст написан на русском для подготовки, с английскими терминами и названиями из презентации. Сами слайды остаются на английском.
\end{titlepage}
\tableofcontents
\newpage
\section{Как пользоваться текстом и распределить время}
Это связный текст выступления, а не только перечень тезисов. Обычные абзацы можно использовать как основу устного объяснения. Серые ремарки адресованы преподавателю; вопросы выделены отдельно, а после них дан разбор, который можно отложить до ответов студентов.

Не нужно стремиться произнести каждую фразу дословно. Главный маршрут: альтернативы и предпочтения; согласованный выбор; полезность как представление; смысл свойств предпочтений; бюджет и выбор на общей картинке. Слайд 37, сепарабельность и квазилинейность на слайдах 46--49, а также подробные доказательства --- резерв. Их включение зависит от темпа группы.

Ориентир для первой пары: 20--25 минут на вступление, организацию курса и кейс; 30--35 минут на предпочтения и наблюдаемый выбор; оставшееся время на полезность и вопросы. Удобная точка остановки --- около слайдов 30--33, но перерыв не обязан совпасть с конкретным слайдом.

Во второй паре отведите примерно 20 минут на непрерывность и монотонность, 10 минут на LNS и две точки насыщения, 20 минут на выпуклость и примеры, 15--20 минут на бюджет и общую картинку. Остаток --- вопросы, итоговый квиз или расширения. Обсуждения и рисунки на доске входят в эти интервалы. Все доказательства приложения за две пары проходить не требуется.

\cue{Перед занятием открыть презентацию в режиме показа и отдельно раздатку. Проверить переход от слайда 24 к доказательству и обратно. На слайде 3 есть ссылка на кейс: если используете её, не нужно второй раз подробно разбирать тот же кейс при линейном проходе слайда 10.}

\section{Введение, организация курса и экономический кейс: слайды 1--12}
\slide{1}{Microeconomics}
Меня зовут Арсений Столяров. У меня PhD in Economics из University of Wisconsin--Madison. В этом модуле мы будем изучать микроэкономику: как описывать решения людей и фирм, какие ограничения на поведение дают наши предпосылки и как использовать эти модели для содержательных объяснений.

У курса будет математический язык, но формулы нужны ради вполне конкретных вопросов. Кто принимает решение? Какие варианты доступны? Как человек или фирма их сравнивает? Что изменится, если поменять цены, доход, технологию или правила? Мне важно, чтобы после решения задачи вы могли ответить на эти вопросы словами.

\cue{Коротко представиться, затем перейти к карте. Не превращать титул в отдельную вводную лекцию.}

\slide{2}{Economics is bigger than this course}
Начнём с масштаба. Экономическая мысль значительно шире того небольшого набора моделей, который мы успеем изучить за семь блоков. На карте показаны несколько крупных традиций. Это не строгая классификация с непересекающимися ящиками и не хронологическое дерево: одни и те же механизмы могут обсуждаться в разных направлениях.

Чтобы названия не были пустыми, возьмём один вопрос: что определяет заработную плату? Классическая политэкономия обращает внимание на производство и распределение между участниками. Марксистская традиция --- на собственность, классовые отношения и присвоение результата. Неоклассическая --- на решения при ограничениях, производительность и относительную редкость. Австрийская --- на распределённое знание и предпринимательское обнаружение возможностей.

Кейнсианские подходы особенно выделяют спрос и занятость, институциональные --- правила и переговорные позиции, поведенческие --- нормы и особенности принятия решений, эволюционные --- изменение технологий и отбор практик. Я не утверждаю, что каждая традиция обладает одним-единственным объяснением или что эти объяснения взаимоисключающие.

Теперь выделим наш участок. Мы в основном будем пользоваться языком неоклассической микроэкономики: потребитель, фирма, риск, выбор во времени и сравнительная статика. Даже внутри этого направления мы берём лишь часть тем. Теорию игр, общее равновесие и дизайн механизмов систематически изучать здесь не будем.

\cue{Показать сначала всю карту; следующим кликом --- рамку курса. Не нужно читать все карточки дословно. Если возникнет вопрос: неоклассическая теория и неолиберализм --- разные понятия, аналитический подход и политико-экономическое направление не следует отождествлять.}

\slide{3}{Why do we need microeconomics?}
Первый ответ: микроэкономика позволяет описывать мир с помощью явных, воспроизводимых моделей. Нам недостаточно сказать «покупатели любят удобство» или «фирма хочет заработать». Нужно указать, что именно выбирается, что ограничено, какие предпосылки приняты и какое наблюдаемое следствие из них получается. Тогда другой исследователь сможет проверить ход рассуждения и сопоставить прогноз с данными.

На практике этот язык нужен для структурных моделей. Например, мы хотим понять, как изменится выбор тарифа после изменения цены. Можно построить статистический прогноз по прошлым данным. Структурная модель дополнительно разделяет предпочтения, ограничения и устройство среды. Мы оцениваем параметры поведения и затем рассматриваем другой тариф или другие правила. Это даёт осмысленный способ строить контрфактические сценарии, хотя их убедительность по-прежнему зависит от идентификации, данных и предпосылок.

Второй ответ: мы начинаем с людей и фирм, а затем разбираемся с агрегированными результатами. После критики Лукаса 1976 года микрофундаменты стали центральной частью значительной доли современной макроэкономики. Если политика меняется, привычная статистическая связь может измениться вместе с поведением людей. Поэтому важно понимать, какие решения породили эту связь. Это не означает, что вся макроэкономика буквально свелась к микроэкономике; это означает, что микроэкономические инструменты вам пригодятся и в центральном банке, и в исследовательском институте.

Третий ответ --- судебные споры. В американских антимонопольных делах и спорах об ущербе экономисты строят модели, оценивают альтернативные сценарии, пишут экспертные заключения и проверяют аргументы противоположной стороны. Одних общих слов о пользе конкуренции там недостаточно: нужно объяснить механизм и показать, какие данные его поддерживают. Отсюда наш пример с Amazon.

\slide{4}{The course at a glance}
В курсе семь лекционных блоков и семь семинаров. Лекции задают язык и логику моделей; семинары нужны, чтобы вы сами применяли определения, строили рисунки и решали задачи.

Итоговая оценка складывается из трёх компонентов: домашние работы дают 20 процентов, проект --- 10 процентов, экзамен --- 70 процентов. Домашних работ будет три, их оцениваемые компоненты имеют одинаковый вес. Проект один: эссе об экономической загадке и короткое представление идеи. Экзамен тоже один, по всему курсу.

Большой вес экзамена означает, что полезно регулярно проверять себя в условиях, похожих на экзаменационные. Поэтому домашние будут устроены специальным образом. Их задача --- не только накопить оценку, но и дать достаточно честную информацию о том, что вы уже можете сделать самостоятельно.

\slide{5}{Each homework has two parts}
В каждой домашней будет две части. Первая --- тренировочный экзамен. По объёму и сложности он будет похож на настоящий экзамен, с учётом уже пройденного материала. Я рекомендую честно выделить 120 минут и попробовать решить его в один заход. Если вы остановились на половине задачи, это ценная информация: значит, именно этот тип рассуждения ещё нужно потренировать.

Семинарист проверит эту часть по-настоящему, чтобы полученный результат что-то говорил о вашей готовности. Но ошибки в ней не будут снижать итоговую оценку за курс. Здесь нет смысла маскировать непонимание красивым решением, которое вы не сможете воспроизвести: тренировочная часть полезна именно своей честностью.

Вторая часть --- обычный оцениваемый набор задач. В среднем задачи там будут сложнее экзаменационных. У вас будет время подумать, обсудить идеи и аккуратно оформить решение. Проверка этой части обычная, с оцениванием результата и аргументации.

Решения домашних нужно писать ручкой на бумаге. Обсуждать подходы можно, но сдаваемое рассуждение должно быть вашим. Использование генеративного ИИ разрешено по правилам курса: укажите инструмент и характер помощи и будьте готовы объяснить любой фрагмент. Умение самостоятельно восстановить логику важнее гладкого текста.

\cue{Не объявлять неоговорённых бонусов за тренировочную часть или дополнительных правил сдачи. Их нет в текущих инструкциях.}

\slide{6}{The project: Why does this happen?}
Проект начинается с наблюдения, которое вас удивляет. Почему подписка устроена именно так? Почему большой размер упаковки иногда невыгоднее маленького? Почему люди выбирают страховку с франшизой? Вопрос должен быть достаточно конкретным, чтобы мы понимали, какое наблюдение вы пытаетесь объяснить.

Дальше нужно выделить принимающего решение, его стимулы и ограничения. Хорошее эссе не просто называет термин. Если написать «это происходит из-за рациональности», мы ещё ничего не объяснили. Нужна причинная цепочка: при таких условиях участник выбирает так-то, поэтому возникает наблюдаемый результат.

Затем полезно спросить: а какое другое объяснение могло бы дать похожее наблюдение? Или что изменилось бы, если бы наш механизм был верен? Один различающий прогноз нередко делает работу гораздо сильнее. Например, если объяснение опирается на неопределённость объёма потребления, реакция людей с предсказуемым и непредсказуемым потреблением должна различаться.

В конце вы сдаёте эссе объёмом примерно 600--900 слов и коротко представляете идею. Нужен источник наблюдения и корректные ссылки. Тему сформулируем после второй лекции, механизм --- после четвёртой; окончательный текст сдаётся после шестой, до седьмой. На последнем семинаре будут презентации и обсуждение. Эти вехи проекта не сдвигаются вместе с домашними.

\slide{7}{The final exam}
Экзамен длится два часа, проходит очно, письменно и индивидуально. В нём будут утверждения true/false с объяснениями, задачи и один качественный вопрос на экономическое рассуждение.

Для true/false недостаточно угадать слово. Если утверждение неверно, хороший контрпример может быть сильнее длинного текста. Если верно, нужно понимать, какие предпосылки обеспечивают вывод. В задачах важны ход решения и интерпретация. В качественном вопросе --- последовательный механизм, а не набор знакомых терминов.

Дату, подробный формат и разрешённые материалы объявим заранее. Сейчас главный практический вывод: решать тренировочную экзаменационную часть домашних стоит так, чтобы получить честный сигнал о готовности.

\slide{8}{Office hours and questions}
Офисные часы --- по запросу. Пишите в Telegram, договоримся о времени. Можно задавать вопросы и по электронной почте, она указана на слайде.

Во время занятия тоже спрашивайте. Особенно полезен вопрос «какую предпосылку мы здесь использовали?». Если формула получилась, а смысл перехода потерялся, лучше остановиться именно в этом месте. Экономическая модель часто ломается не из-за сложного вычисления, а из-за незамеченного изменения условий задачи.

\slide{9}{Seven lecture blocks: the topic map}
В первом блоке мы строим язык предпочтений, выбора и полезности. На первом семинаре к этому добавятся бинарные отношения, кривые безразличия и бюджетные множества. Во втором блоке систематически решаем задачу потребителя и получаем маршаллианский спрос, косвенную полезность и тождество Роя.

Третий блок --- дуальность, разложение Слуцкого и выявленные предпочтения. Четвёртый --- выбор в условиях риска. Пятый --- фирма: технология, прибыль и издержки. Шестой --- монотонная сравнительная статика, возрастающие разности, супермодулярность и Топкис. Седьмой --- межвременной выбор и финансовые приложения.

Домашние теперь выдаются после третьего, пятого и седьмого блоков. Это обновлённый график. Сроки сдачи будут объявлены отдельно; на этом слайде отмечены моменты выдачи. Проект идёт параллельно по своим вехам.

Обратите внимание на общую структуру: сначала учимся описывать желания и возможности, потом решать конкретные задачи, затем изучаем реакцию решений на изменение среды. На сегодняшней лекции не нужно заранее освоить все методы оптимизации, которые понадобятся дальше.

\slide{10}{Amazon and off-platform prices}
В 2023 году FTC и штаты предъявили Amazon антимонопольные обвинения. Один из обсуждаемых в жалобе механизмов --- санкции против продавцов, которые предлагают более низкие цены на конкурирующих площадках. Здесь мы разбираем экономическую логику обвинения, а не утверждаем, что суд уже установил нарушение.

Возьмём условного продавца. Себестоимость единицы --- 60 долларов. Условная плата площадке составляет 20 долларов на Amazon и 10 долларов на eBay. Если продавать за 100 на первой площадке и за 90 на второй, прибыль на единицу в обоих случаях равна 20. Все эти числа учебные: это не реальные тарифы площадок.

Теперь представим, что более низкая цена на другой площадке ведёт к потере видимости или продаж на Amazon. У продавца появляется дополнительный стимул держать внешнюю цену высокой. Тогда меньшая комиссия конкурента не обязательно превратится в более низкую цену для покупателя. Так ограничение на одной площадке может повлиять на конкуренцию между площадками.

\ask{Кто выигрывает и кто проигрывает? Обязательно ли рост внешней цены --- единственная возможная реакция продавца?}
\answer{Продавец может поднять внешнюю цену, изменить ассортимент, уйти с одной площадки или принять потерю части продаж. Покупатели могут столкнуться с меньшим разбросом цен и более высокими ценами; для точного вывода нужно знать спрос и варианты переключения. Нельзя оценить общий эффект, посмотрев только на одну комиссию.}

Можно ли привести аргумент в защиту ограничений? Например, площадка может утверждать, что инвестирует в поиск, доверие и привлечение покупателей, а покупка затем уходит в другой канал. Это гипотеза о free riding. Её тоже надо проверять: какие инвестиции действительно защищаются, существует ли такая проблема, есть ли менее ограничительные способы её решить?

Экономист в таком споре должен связать правило, стимулы продавца, поведение покупателя и контрфактический мир без правила. Понадобятся данные о ценах, видимости, комиссиях, переключении и ассортименте. Именно ради таких рассуждений нам нужен язык микроэкономики.

\cue{Это FTC and states v. Amazon, а не иск Amazon к eBay. В текущей подготовке официальный сайт FTC был недоступен через исходящий прокси; актуальный процессуальный статус не проверен. В речи достаточно точно сказать «в жалобе 2023 года утверждалось». Ссылка на источник есть в презентации.}

\slide{11}{Today: from primitives to useful assumptions}
Сегодня у нас две связанные задачи. В первой части мы описываем выбор: сначала через предпочтения, затем через правило выбора из меню. Разберём, что означает рациональность, какие ограничения даёт согласованность выбора и когда полезность позволяет представить эту задачу числами.

Во второй части посмотрим на свойства предпочтений: непрерывность, монотонность, локальную ненасыщаемость и выпуклость. Это не просто список технических слов. Каждая предпосылка что-то исключает и что-то позволяет доказать. Если останется время, обсудим сепарабельность и квазилинейность --- два распространённых способа упростить модели.

Завершим бюджетным множеством и совместной картинкой желаний и возможностей. Это подготовит нас к следующей лекции о максимизации полезности.

\slide{12}{To model a choice, we need to know...}
Чтобы моделировать выбор, нам нужно понять две вещи.

Что люди хотят?

И что они могут?

Первое мы опишем через предпочтения. Второе --- через множество доступных альтернатив. Только соединяя эти два объекта, можно обсуждать выбор. Любимый товар может быть недоступен, а купленный товар не обязан быть самым любимым из всех мыслимых.

Начнём с того, чего люди хотят.

\cue{Три смысловых шага на слайде раскрываются отдельно. Дать паузу между «хотят» и «могут».}

\section{Предпочтения и согласованность выбора: слайды 13--26}
\slide{13}{What is the primitive?}
Прежде чем сравнивать альтернативы, определим, что вообще считается альтернативой. Пусть $X$ --- множество полностью описанных вариантов. Это может быть конечный список тарифов: basic, standard, pro. Может быть непрерывное пространство наборов $X=\R_+^L$, где каждая координата означает количество отдельного блага.

Возможны и смешанные, дискретно-непрерывные альтернативы. Например, выбираем чай или кофе и одновременно количество молока. Упрощённо запишем
\[
X=\{\text{tea},\text{coffee}\}\times\R_+.
\]
Первая компонента дискретна, вторая непрерывна. В реальном стакане количество молока ограничено объёмом; при необходимости это легко отразить в множестве $X$.

Важный момент: одинаковое количество запросов к ИИ-сервису с разными гарантиями приватности может быть разными альтернативами. То же относится к качеству, сроку, условиям доступа и информации. Если мы забыли включить существенную характеристику, дальнейшая математика может быть правильной, а модель --- плохо описывать реальный выбор.

\slide{14}{One relation is enough to start}
Введём одно бинарное отношение на $X$: слабое предпочтение $\succeq$. Запись $x\succeq y$ читается как «$x$ не хуже $y$», или «$x$ по крайней мере так же хорош, как $y$».

Формально отношение --- подмножество $X\times X$: для каких-то пар сравнение верно, для каких-то нет. Пока мы не говорим, какие свойства у этого отношения. Оно ещё не обязано сравнивать все пары или быть согласованным по цепочке.

Символ похож на числовое $\ge$, но это не числовое сравнение. Мы пока не присваивали альтернативам оценки. И это не описание покупки: человек может предпочитать недоступный вариант. Это отдельный объект модели --- способ сравнивать пары альтернатив.

В начальных изложениях иногда дополнительно вводят безразличие и строгое предпочтение. Давайте проверим, нужны ли нам для этого новые независимые исходные объекты.

\slide{15}{Can we define indifference using only the weak relation?}
Мы уже умеем сказать «$x$ не хуже $y$». Как сказать «они одинаково хороши»?

\ask{Как записать безразличие, используя только $\succeq$?}
\cue{Сначала выслушать ответ, затем открыть формулу.}
\answer{Нужно потребовать оба слабых сравнения:
\[
x\sim y\quad\Longleftrightarrow\quad x\succeq y\ \text{и}\ y\succeq x.
\]}

Обратите внимание на отличие безразличия от несравнимости. При безразличии оба сравнения верны. При несравнимости не верно ни одно: отношение просто не ранжирует эту пару. Если я говорю «не знаю, как сопоставить эти варианты», это не то же самое, что «для меня они равноценны».

Пока мы ещё не приняли полноту, поэтому несравнимость логически возможна. Безразличие мы определили, не добавляя второго независимого отношения.

\slide{16}{Can we define strict preference using only the weak relation?}
Теперь как отличить «не хуже» от «строго лучше»? Одного $x\succeq y$ недостаточно: возможно, $y$ тоже не хуже $x$.

\ask{Какое условие нужно добавить?}
\answer{Нужно исключить обратное слабое предпочтение:
\[
x\succ y\quad\Longleftrightarrow\quad
x\succeq y\ \text{и не}\ y\succeq x.
\]}

То есть из одного слабого отношения получаются и строгое предпочтение, и безразличие. Новые символы удобны, но дополнительными примитивами модели они не являются.

\ask{Если неверно $x\succ y$, обязательно ли верно $y\succ x$?}
\answer{Нет. Даже если все пары сравнимы, возможно $x\sim y$. А без полноты возможна ещё и несравнимость. Отрицание строгого предпочтения нельзя автоматически заменить обратным строгим предпочтением.}

\slide{17}{What does rationality mean in a model?}
Экономисты любят предполагать рациональные предпочтения. Но слово «рациональный» в повседневной речи несёт слишком много смыслов. Что действительно нужно нашей модели?

\ask{Какие требования согласованности вы бы предъявили к ранжированию альтернатив?}

Первое --- полнота. Для любых $x,y\in X$ верно $x\succeq y$ или $y\succeq x$, возможно оба сравнения. Никакая пара не остаётся несравнимой. Полнота допускает безразличие и не говорит, что сравнение легко сделать. Если подставить $x=y$, получим рефлексивность: $x\succeq x$.

Второе --- транзитивность. Если $x$ не хуже $y$, а $y$ не хуже $z$, то $x$ не хуже $z$:
\[
x\succeq y,\quad y\succeq z\quad\Rightarrow\quad x\succeq z.
\]

В этом курсе рациональность предпочтений означает именно полноту и транзитивность. Из определения не следуют эгоизм, идеальная осведомлённость или способность мгновенно решать сложные вычислительные задачи. Альтруистические предпочтения тоже могут быть рациональными. Это термин о структуре ранжирования, а не похвала человеку.

\slide{18}{Same customer, two restaurant menus}
Представьте одного посетителя и два меню. В первом есть стейк, утиное конфи с картофельным гратеном и котлета с пюре. Он выбирает стейк. Во втором есть тот же стейк, то же конфи и дегустационное меню из семи блюд, например с равиоли из омара и олениной с трюфельным соусом. Теперь посетитель выбирает конфи.

\ask{Как можно объяснить такую смену выбора?}
\cue{Не открывать сразу последние допущения. Дать студентам перечислить версии.}

Можно предположить разное качество, цены, ожидания от ресторана, настроение или контекст. Можно увидеть эффект компромисса: в первом меню стейк выглядит средней по сложности опцией, во втором такой опцией оказывается конфи. Если человек выбирает относительно других пунктов меню, описание через одно неизменное ранжирование самих блюд может оказаться недостаточным.

Теперь уточним: общие блюда, цены и информация одинаковы; остальные важные обстоятельства фиксированы; каждый выбор показывает единственного фаворита. Тогда первое меню говорит «стейк строго лучше конфи», а второе --- обратное. Одно фиксированное ранжирование этого объяснить не может.

Здесь важно не сделать слишком быстрый вывод. Две покупки сами по себе не доказывают нарушение трёхзвенной транзитивности. Если стейк и конфи равноценны, а правило разрешения ничьей меняется, обе покупки могут быть совместимы с рациональными предпочтениями. Вопрос прежде всего о согласованности выбора из разных меню. Сейчас мы сформулируем соответствующую аксиому.

\slide{19}{Your turn: three relations}
Проверим определения на трёх отношениях. В первом сравниваем суммы координат: $x\succeq y$ тогда и только тогда, когда $x_1+x_2\ge y_1+y_2$. Во втором сравниваем трёхмерные векторы по правилу «не меньше хотя бы по двум из трёх координат». В третьем требуем не меньшего количества каждого из двух благ.

\ask{Какие отношения полны? Какие транзитивны? Для нарушения свойства нужен конкретный пример.}
\cue{Дать время подумать и собрать ответы; числовую подсказку ко второму отношению показывать после первых попыток.}

Первое отношение полно и транзитивно: мы просто сравниваем вещественные числа --- суммы. Равенство сумм даёт безразличие.

Второе полно. Если в одном направлении не нашлось хотя бы двух слабых координатных сравнений, в другом направлении они обязательно найдутся. Но транзитивности нет. Возьмём
\[
a=(3,2,1),\quad b=(2,1,3),\quad c=(1,3,2).
\]
$a$ выигрывает у $b$ по первым двум координатам, $b$ у $c$ --- по первой и третьей, а $c$ у $a$ --- по второй и третьей. Получился строгий цикл. Правило выглядит разумным по каждой отдельной паре, но глобально несогласованно.

Третье отношение транзитивно: координатные неравенства складываются по цепочке. Полноты нет: $(2,1)$ и $(1,2)$ несравнимы. Здесь нельзя путать отсутствие сравнения с безразличием.

\slide{20}{Condorcet's paradox: rational voters, cyclic majority}
Только что мы получили цикл из правила сравнения по большинству координат. Теперь дадим этим координатам содержательный смысл: пусть каждая координата представляет одного избирателя. Получится парадокс Кондорсе.

Есть три альтернативы $a,b,c$ и три человека. Первый ранжирует их как $a\succ_1 b\succ_1 c$. Второй --- как $c\succ_2 a\succ_2 b$. Третий --- как $b\succ_3 c\succ_3 a$. У каждого человека обычное строгое, полное и транзитивное ранжирование. Ни один из них не попадает в цикл.

Сравним альтернативы попарным большинством. Между $a$ и $b$ первый и второй выбирают $a$, третий --- $b$. Значит, $a$ побеждает $b$ со счётом два к одному. Между $b$ и $c$ первый и третий поддерживают $b$: снова два к одному. Между $c$ и $a$ второй и третий поддерживают $c$.

Итак, коллективное отношение выглядит так:
\[
a\succ_M b,\qquad b\succ_M c,\qquad c\succ_M a.
\]
Индивидуальная рациональность не гарантировала транзитивности коллективного ранжирования. Здесь нет победителя Кондорсе --- альтернативы, которая побеждает каждую другую в парном голосовании большинством.

\ask{Что произойдёт, если сначала голосовать между $a$ и $b$, а потом между победителем и $c$? А если начать с $b$ и $c$?}
\answer{В первом случае итогом будет $c$, во втором --- $a$. Порядок постановки вопросов может влиять на результат, хотя предпочтения избирателей не менялись.}

Это не утверждение, что большинство всегда циклично или что сами избиратели ведут себя нетранзитивно. Это конкретный пример того, что свойство индивидуальных отношений может потеряться при агрегировании. Поэтому в модели коллективного выбора недостаточно просто написать «все рациональны».

\cue{На рисунке стрелка идёт от победителя к проигравшему. Открывать три ребра по одному. В основной маршрут достаточно 3--5 минут; вопрос о порядке голосования --- дополнительная минута по реакции группы.}

\slide{21}{We see choices, not the whole preference relation}
Теория предпочтений сама по себе полезна, но в реальности мы не видим полное ранжирование каждого человека. Обычно у нас есть покупки, клики или решения из нескольких конкретных меню. Поэтому полезно спросить, что о предпочтениях можно узнать из выбора.

Допустим, из меню $\{A,B,C\}$ выбирается $A$, а из $\{B,C\}$ --- $B$. Это совместимо с неизменным порядком $A\succ B\succ C$: вторая покупка другая, потому что исчез лучший доступный вариант. Изменение поведения не доказывает изменение вкусов.

Но два наблюдения не дают автоматически всех знаний о человеке. Мы должны знать доступные варианты, их характеристики и контекст. Ещё важно, что именно означает слово «выбирается»: наблюдаем ли мы единственный реализованный выбор или всё множество лучших альтернатив? Сейчас это различие станет формальным.

\ask{Чего нам не хватает, если в данных записано только «человек купил тариф standard»?}
\answer{Нужно знать, какие другие тарифы были доступны, их условия и цены, информацию покупателя и то, как он разрешает безразличие. Иначе даже сильное утверждение о сравнении двух тарифов может оказаться необоснованным.}

\slide{22}{Choice as the primitive}
Можно начать модель не с предпочтений, а с правила выбора. Пусть $\mathcal D$ --- множество всех непустых меню, то есть всех непустых подмножеств $X$. Обозначение $2^X$ означает множество подмножеств, а не возведение какого-то количества товаров в степень.

Правило $C$ каждому меню $A$ сопоставляет непустое подмножество выбранных вариантов:
\[
C:\mathcal D\to2^X,\qquad\varnothing\ne C(A)\subseteq A.
\]
То есть правило не выбирает то, чего нет в меню, и на непустом меню выбирает хотя бы что-то. Результат может быть многозначным. Если чай и кофе одинаково хороши, оба входят в $C(\{\text{чай},\text{кофе}\})$, хотя в конкретный момент человек может купить только один напиток.

Если у нас уже есть предпочтения, естественно определить $C(A)$ как все элементы $A$, которые не хуже каждого другого элемента меню:
\[
C(A;\succeq)=\{x\in A:x\succeq y\ \text{для всех }y\in A\}.
\]
Так мы идём от ранжирования к выбору. Но можно пойти в другую сторону: задать правило выбора и проверить, ведёт ли оно себя так, как будто за ним стоит одно рациональное ранжирование.

На бесконечном меню непустота множества лучших вариантов не гарантирована. Если выбирать число из $(0,1)$ и больше всегда лучше, максимального числа нет. Поэтому требование непустоты не стоит незаметно смешивать с определением предпочтений.

\slide{23}{Consistency across menus: WARP}
Нам нужно условие согласованности выбора между меню. Оно называется Weak Axiom of Revealed Preference, или WARP.

Возьмём два меню $A$ и $B$ и две альтернативы $x,y$, которые доступны в обоих. Если $x$ выбран в $A$, а $y$ выбран в $B$, то аксиома требует, чтобы $x$ тоже был выбран в $B$, а $y$ --- в $A$:
\[
\begin{gathered}
x,y\in A\cap B,\qquad x\in C(A),\quad y\in C(B)\\
\Rightarrow\quad x\in C(B),\quad y\in C(A).
\end{gathered}
\]

Почему это естественно для выбора лучших вариантов? В первом меню $x$ не хуже $y$, ведь $y$ тоже доступен. Во втором $y$ не хуже $x$. При одном рациональном ранжировании они должны быть безразличны. Если один из них лучший в соответствующем меню, второй тоже должен быть среди лучших.

Это не запрет менять покупку. Когда любимый товар исчезает, покупка может измениться без всякой проблемы. Аксиома относится к двум общим доступным альтернативам и к целому множеству лучших вариантов, включая ничьи. Она не говорит, что любое различие двух наблюдённых чеков --- нарушение рациональности.

\cue{Пройти кванторы словами. В приложениях к бюджетным данным формулировки могут выглядеть иначе; здесь используется версия для полного соответствия выбора.}

\slide{24}{A two-menu inconsistency}
Рассмотрим $A=\{a,b,c\}$ и $B=\{a,b,d\}$. Пусть мы знаем именно полные выбранные множества: $C(A)=\{a\}$ и $C(B)=\{b\}$. И $a$, и $b$ доступны в обоих меню.

\ask{Может ли одно полное транзитивное отношение породить такие выбранные множества?}
\answer{Нет. Из первого следует, что $a$ строго лучше $b$: $a$ лучший, а $b$ не входит в множество лучших. Из второго следует обратное строгое сравнение. Эти требования несовместимы.}

Если бы мы просто увидели одну покупку $a$ и одну покупку $b$, не зная, какие ещё варианты были равноценными, вывод был бы слабее. Нельзя незаметно заменить данные об одном реализованном действии данными о всём $C(A)$. Мы будем постоянно возвращаться к вопросу о том, что именно наблюдается.

\slide{25}{Preferences and choice: an equivalence}
Теперь сформулируем связь точно. Пусть $X$ конечно, а $C$ непусто на каждом непустом подмножестве $X$. Тогда WARP выполняется тогда и только тогда, когда существует полное транзитивное отношение, для которого $C$ выбирает все лучшие элементы каждого меню.

В одну сторону результат интуитивен: рациональное ранжирование задаёт согласованный выбор. В другую он говорит, что достаточно согласованное правило выбора можно рационализировать предпочтениями. Для восстановления ранжирования смотрим на парные меню:
\[
x\succeq y\quad\Longleftrightarrow\quad x\in C(\{x,y\}).
\]

Здесь важны все условия. Конечность снимает трудности существования лучших вариантов. Непустота позволяет прочитать хотя бы одно сравнение из каждой пары. Полный набор меню даёт возможность сопоставлять пары и тройки. Нельзя заменить «каждое меню» на «несколько меню в моей выборке» и сохранить тот же вывод без дополнительных рассуждений.

\cue{На слайде есть переходы к обоим направлениям доказательства. Полный текст находится в разделах к слайдам 55--56 ниже. Если времени мало, оставить теорему и интерпретацию, не проходить доказательство целиком.}

\slide{26}{What choice data do not automatically reveal}
У теоремы есть два важных эмпирических ограничения. Первое: обычно мы наблюдаем один реализованный выбор, а не всё множество лучших вариантов. Если мы никогда не узнаём, что точно не входит в $C(A)$, трудно исключить даже универсальное безразличие. Человек мог считать все варианты одинаково хорошими и каждый раз разрешать ничью по-разному.

Второе: наблюдаются только некоторые меню. Цикл может спрятаться в неполных данных. Если мы увидели три парных выбора, образующих «камень, ножницы, бумага», разные парные меню имеют лишь по одной общей альтернативе. На наблюдаемом домене WARP может не поймать нарушение, хотя единого рационального ранжирования, порождающего эти полные выбранные множества, не существует.

Отсюда практический вывод для данных о рекомендациях, тарифах и покупках: содержание вывода зависит от того, какие меню наблюдались, что известно о ничьих и какие дополнительные предпосылки мы приняли. Наличие числовой метрики или большого лога действий само по себе не решает проблему идентификации предпочтений.

\cue{Подробный пример неполных меню дан к слайду 57. Переходя дальше, подчеркнуть: теперь мы не добавляем новые вкусы, а ищем удобный числовой язык для уже заданного ранжирования.}

\section{Полезность как представление: слайды 27--33}
\slide{27}{Utility is a representation}
Функция полезности --- вещественная функция $u:X\to\R$. Она представляет предпочтения, если для любых двух альтернатив
\[
x\succeq y\quad\Longleftrightarrow\quad u(x)\ge u(y).
\]
Нужны обе стороны «тогда и только тогда». Не только более высоко оценённые варианты должны быть не хуже; вся структура сравнений, включая безразличие, должна совпасть.

Из определения сразу следует: $x\succ y$ тогда и только тогда, когда $u(x)>u(y)$, а $x\sim y$ тогда и только тогда, когда значения равны. Когда представление есть, поиск лучших доступных вариантов можно записать как оптимизацию:
\[
C(A;\succeq)=\operatorname*{arg\,max}_{x\in A}u(x).
\]

В этом и удобство полезности: мы умеем работать с функциями и задачами оптимизации. Но числовая запись ещё не доказывает, что максимум достигается. На бесконечном меню argmax может быть пустым. Сначала определили представление, затем отдельно обсуждаем его существование и существование решения.

\slide{28}{Representation on a finite set}
На конечном множестве любое полное транзитивное отношение имеет числовое представление. Построение очень простое. Сначала объединим все безразличные альтернативы в классы. Транзитивность делает такое объединение согласованным: если $x$ безразличен $y$, а $y$ безразличен $z$, то и $x$ безразличен $z$.

Дальше упорядочим классы от худшего к лучшему и дадим им возрастающие числовые метки. Например, $C$ и $D$ получают ноль, $B$ --- единицу, $A$ и $E$ --- двойку. Внутри класса одинаковое значение, между классами правильный порядок.

Здесь нам не пришлось измерять какое-либо вещество под названием «полезность». Мы просто закодировали ранжирование. Можно было выбрать другие возрастающие метки и получить то же поведение. Именно это означает порядковый, ordinal, характер представления.

\slide{29}{Ordinal, not a happiness meter}
Посмотрите на две колонки. В первой альтернативам $A,B,C$ присвоены числа $3,2,1$. Во второй --- $100,7,-20$. Обе колонки описывают одно и то же: $A\succ B\succ C$.

Из первой колонки нельзя заключить, что $A$ в три раза лучше $C$. Нельзя также интерпретировать одинаковые разности между соседними значениями как одинаковое количество счастья. Мы пока не задали ни единицы измерения, ни экономического смысла нуля, ни количественного смысла расстояний.

Покупатель ведёт себя так, как будто максимизирует подходящую функцию. Это удобное представление его выбора, а не обязательная гипотеза о том, что он вычисляет числа в голове или что существует прибор, измеряющий именно эти значения.

\slide{30}{Increasing transformations preserve the ranking}
Если $u$ представляет предпочтения, то любое строго возрастающее преобразование $v=f\circ u$ представляет те же предпочтения. Причина в сохранении порядка:
\[
u(x)\ge u(y)\quad\Longleftrightarrow\quad f(u(x))\ge f(u(y)).
\]
Можно прибавить пять, умножить на положительное число, применить экспоненту. Возведение в квадрат подходит на неотрицательном диапазоне, но не в общем случае: например, $-1>-2$, а квадраты упорядочены наоборот.

\ask{Почему недостаточно потребовать, чтобы $f$ была просто неубывающей?}
\cue{Ответа на самом слайде больше нет. Сначала дать студентам найти контрпример.}
\answer{Неубывающая функция может быть плоской. Если $u(A)=2>1=u(B)$, а $f(t)=0$, новое представление делает $A$ и $B$ безразличными. Оно изменило исходное отношение. Нужна строгая монотонность на реально достигаемом диапазоне $u(X)$; вне этого диапазона поведение $f$ неважно.}

Связь с ML здесь ограниченная, но полезная. Монотонное преобразование уже определённого детерминированного целевого показателя сохраняет argmax. Однако преобразование каждой награды перед суммированием или преобразование до взятия ожидания может изменить оптимум. Мы вернёмся к этому различию, когда будем обсуждать риск.

\slide{31}{No interpersonal comparisons for free}
Пусть Яромир и Зоя оба предпочитают $x$ набору $y$. У каждого есть своё порядковое представление. Полезность Яромира можно независимо умножить на миллион, не изменив ни одного его сравнения. Полезность Зои можно независимо сжать до очень маленького диапазона, тоже ничего не изменив в её предпочтениях.

Но если теперь сложить эти функции, веса людей в сумме фактически поменяются. Значит, одно только наличие двух вещественных представлений не даёт общей единицы благополучия и не оправдывает произвольную сумму как общественный критерий.

Это не запрет на анализ благосостояния. Это напоминание, что для него нужны дополнительные нормативные или межличностные предпосылки. Их нельзя извлечь из чисел, которые первоначально были лишь метками порядка.

\slide{32}{Does every ranking have a utility representation?}
Мы доказали существование представления на конечном множестве. А можно ли таким способом представить любые рациональные предпочтения на бесконечном пространстве?

\cue{Задать вопрос до раскрытия примера. Важно, что речь о полном транзитивном отношении, а не о заведомо циклическом.}

Рассмотрим лексикографические предпочтения на $[0,1]^2$. Сначала сравниваем первую координату и только при её равенстве переходим ко второй:
\[
x\succeq_{\mathrm{lex}} y
\iff [x_1>y_1]\ \text{или}\ [x_1=y_1\ \text{и}\ x_2\ge y_2].
\]
Это похоже на упорядочение слов по алфавиту. Любое, сколь угодно маленькое преимущество по первой координате важнее любой разницы по второй. Отношение полно и транзитивно, но вещественной функции полезности, представляющей его на всём квадрате, нет.

Интуитивно нельзя просто поставить перед первой координатой очень большой конечный коэффициент. Как бы велик он ни был, первая координата может отличаться ещё меньше, а разница по второй его перевесит. Это объясняет, почему обычная взвешенная сумма не работает; доказательство невозможности вообще любой вещественной функции сильнее и использует счётность рациональных чисел.

На конечном списке лексикографический порядок представить можно. Проблема появляется из-за непрерывного бесконечного множества возможных первых координат, а не из-за «иррациональности» самого правила.

\slide{33}{Checkpoint}
Перед следующим шагом проверим четыре вещи. Нарушает ли безразличие полноту? Нет: полнота разрешает оба слабых сравнения. Может ли полезность представлять нетранзитивные предпочтения? Нет: порядок вещественных чисел транзитивен, и представление переносит это свойство.

Доказывает ли другая покупка изменение предпочтений? Нет: могли измениться меню, цены, информация или разрешение ничьей. И наконец: что изменилось при переходе от конечного множества к бесконечному?

Здесь два отдельных вопроса. Во-первых, рациональность больше не гарантирует вещественное представление: пример --- лексикографический порядок. Во-вторых, даже существующая полезность не гарантирует достижения максимума: $u(x)=x$ на $(0,1)$ имеет представление, но не имеет максимизатора.

Эти две трудности требуют разных условий. Следующий шаг --- условия, при которых существует хорошее представление. Позднее отдельно разберём условия существования оптимума.

\cue{Если наступил перерыв, остановиться здесь. Если часть тем первой пары заняла больше времени, допустимо сделать перерыв раньше и начать вторую с этого краткого восстановления логики.}

\section{Непрерывность, монотонность и насыщение: слайды 34--39}
\slide{34}{Continuity: comparisons survive limits}
Теперь альтернативы --- точки евклидова пространства. Мы можем говорить о близости и пределе последовательностей. Непрерывность предпочтений означает, что слабые сравнения сохраняются при переходе к пределу:
\[
\begin{gathered}
x^n\to x,\quad y^n\to y,\quad x^n\succeq y^n\ \text{при всех }n\\
\Rightarrow\quad x\succeq y.
\end{gathered}
\]

Идея в том, что сравнение не должно внезапно переворачиваться в пределе. Лексикографический порядок этот тест не проходит. Для всех $n$ набор $(1/n,0)$ строго лучше $(0,1)$, потому что первая координата положительна. Но его предел $(0,0)$ строго хуже $(0,1)$: первые координаты равны, а вторая меньше.

Стандартная теорема говорит: полные, транзитивные и непрерывные предпочтения на $\R_+^L$ имеют непрерывное представление полезностью. Здесь существование некоторой функции и существование непрерывной функции --- не одно и то же. Более того, если предпочтения непрерывны, не всякая случайно выбранная числовая запись обязана быть непрерывной: строго возрастающая функция с разрывом может изменить свойства записи, сохранив порядок.

Общую теорему целиком сейчас не доказываем. После монотонности покажем более наглядное построение при дополнительном условии. Главное --- понять, какую проблему решает непрерывность и почему рациональности самой по себе было мало.

\slide{35}{Indifference curves: equally good bundles}
Когда есть полезность, удобно изображать множества одинаково хороших наборов. Для фиксированного $\bar x$ его множество безразличия имеет вид
\[
I(\bar x)=\{x:x\sim\bar x\}=\{x:u(x)=u(\bar x)\}.
\]
В пространстве двух благ мы часто получаем кривую безразличия, indifference curve. На рисунке взята конкретная функция $u(x)=\sqrt{x_1x_2}$. Одно и то же значение полезности означает, что произведение координат постоянно.

Каждая такая кривая объединяет равноценные наборы. После строго возрастающего преобразования полезности числовые метки поменяются, а сами кривые останутся теми же. Например, переход от $\sqrt{x_1x_2}$ к $x_1x_2$ меняет подписи уровней, но не ранжирование.

Эта картинка описывает желания. Пока мы не добавили бюджет, она не говорит, что доступно. Разные классы безразличия при рациональных предпочтениях не могут иметь общую точку: иначе транзитивность безразличия объединила бы их в один класс. Сейчас будем связывать свойства ранжирования с геометрией этих множеств.

\slide{36}{Monotonicity: fix the notation first}
Нужно аккуратно договориться об обозначениях. $x\ge y$ означает не меньше по каждой координате. $x\gg y$ означает строго больше по каждой координате. Это координатные сравнения количеств, не сами предпочтения.

В нашей терминологии monotonicity означает: если выросли все количества, набор стал строго лучше. Strong monotonicity сильнее: если ни одно количество не уменьшилось и хотя бы одно выросло, набор стал строго лучше:
\[
\begin{aligned}
x\gg y&\Rightarrow x\succ y,\\
x\ge y,\ x\ne y&\Rightarrow x\succ y\quad\text{при сильной монотонности}.
\end{aligned}
\]

Названия в разных учебниках отличаются, поэтому надёжнее смотреть на формулы. Обе предпосылки дополняют рациональность, а не следуют из неё. Полнота и транзитивность не обещают, что больше всегда лучше.

Это также содержательное условие о выборе благ. Если координата --- шум или загрязнение, монотонность по ней может быть неуместна. Иногда полезно переопределить товар как тишину или качество воздуха. Иногда лучше прямо моделировать насыщение или антиблаго.

\slide{37}{A constructive picture: the diagonal bundle}
\reserve
Можно наглядно построить полезность, если к рациональности и непрерывности добавить нашу монотонность. Рассмотрим диагональные наборы $\alpha e$, где $e=(1,\ldots,1)$. На диагонали всех благ поровну.

Для произвольного $x$ найдём диагональный набор, который ему безразличен. Маленький диагональный набор будет не лучше $x$, достаточно большой --- не хуже. Непрерывность позволяет найти точку безразличия между ними. А два разных диагональных набора строго ранжируются по монотонности, поэтому такой коэффициент единственен.

Определим полезность $u(x)$ как этот коэффициент $\alpha$. Если $x$ соответствует более высокому диагональному набору, чем $y$, то и сам $x$ лучше $y$. Значит, коэффициент воспроизводит ранжирование.

\pictureof{37}

В примере на рисунке $(1,4)$ безразличен $(2,2)$ при $u=\sqrt{x_1x_2}$, поэтому его диагональная метка равна двум. Сам рисунок иллюстрирует построение, но не заменяет доказательство для всех предпочтений. Полное рассуждение, включая непрерывность построенной функции, есть в приложении.

Заметьте: общая теорема о непрерывном представлении не требует монотонности. Она нужна именно этому удобному способу построения.

\slide{38}{Local nonsatiation: a better bundle nearby}
Представьте кофе и пирожное. Может ли быть такой набор, что любое достаточно маленькое изменение уже не делает вас лучше? Вполне можно представить идеальное для вас сочетание и количество. Тогда предпосылка, которую сейчас введём, в этой точке не выполнится.

Локальная ненасыщаемость, local nonsatiation или LNS, говорит противоположное: около любого набора есть строго лучший, как бы маленькую окрестность мы ни взяли:
\[
\forall x\in X\ \forall\varepsilon>0\ \exists y\in X:
\|y-x\|<\varepsilon,\qquad y\succ x.
\]

Не нужно, чтобы любое увеличение любого блага помогало. Достаточно хотя бы одного направления улучшения поблизости. Например, $u(x)=x_1$ локально ненасыщаема: можно немного увеличить первое благо. Второе вообще не влияет на полезность, поэтому сильной монотонности нет. А $u=x_1-x_2$ показывает, что LNS не требует даже нашей монотонности по всем благам: второе здесь антиблаго.

На полном неотрицательном ортанте сильная монотонность влечёт монотонность, а та --- LNS. Обратные импликации в общем случае неверны.

Зачем условие полезно? При положительных ценах и обычном линейном бюджете оно обеспечивает исчерпание бюджета в существующем оптимуме. Если деньги остались, достаточно близкие наборы ещё доступны, а LNS даёт среди них строго лучший. Но существование и единственность решения LNS сама по себе не гарантирует.

\slide{39}{Two bliss points}
На следующей картинке два набора, $a$ и $b$, являются равноценными точками насыщения. Приближение к любой из них улучшает положение, а в самих точках строго лучшего набора уже нет.

\pictureof{39}
\ask{Могут ли такие предпочтения быть рациональными и непрерывными? Выполняется ли LNS? А выпуклость?}

\answer{Рациональность и непрерывность вполне возможны. Например,
\[
u(x)=-\min\{\|x-a\|^2,\|x-b\|^2\},\quad a=(2,2),\ b=(5,4).
\]
Это непрерывная вещественная функция, значит, она задаёт полное транзитивное отношение. Обе вершины имеют полезность ноль. LNS в них нарушена: строго лучшего набора нет ни рядом, ни вообще.}

Выпуклости тоже нет. Среднее двух лучших наборов --- $(3.5,3)$ --- хуже каждого из них: его полезность $-13/4$, а в вершинах ноль. Раз мы смешали два одинаково хороших набора и получили худший, вкус к смешению нарушен.

Эта картинка полезна сразу для нескольких различий. Рациональность не означает ненасыщаемость. Непрерывность не означает единственную точку насыщения. А геометрия предпочтений может существенно отличаться от привычных нисходящих кривых. Сейчас определим условие, которое ограничивает именно отношение к смешению.

\section{Выпуклость и три стандартных примера: слайды 40--45}
\slide{40}{Convex preferences: mixing bundles}
Пусть пространство наборов выпукло. Предпочтения выпуклы, если смесь двух наборов, каждый из которых не хуже некоторого $z$, тоже не хуже $z$:
\[
x\succeq z,\ y\succeq z\Rightarrow\lambda x+(1-\lambda)y\succeq z,
\qquad\lambda\in[0,1].
\]
Геометрически все верхние контурные множества должны быть выпуклы. Если два набора безразличны, их среднее не хуже каждого из них, но строгого выигрыша обычная выпуклость не обещает.

Строгая выпуклость усиливает вывод: для различных $x,y$ и внутренней смеси $0<\lambda<1$ получаем строгое предпочтение смеси общему уровню $z$. В частности, смешение двух разных безразличных наборов строго помогает. Условие различия необходимо: смесь одного набора с самим собой не может стать лучше него.

Мы смешиваем количества благ, а не разыгрываем лотерею. Например, половина $(4,0)$ плюс половина $(0,4)$ --- набор $(2,2)$. Это не случайное получение одного из крайних наборов. Теперь разберёмся, как данное экономическое свойство связано с формой представляющей функции.

\slide{41}{Convex functions and convex preferences}
Сначала отдельно определим выпуклую функцию. На выпуклой области $u$ выпукла, если
\[
u(\lambda x+(1-\lambda)y)\le\lambda u(x)+(1-\lambda)u(y).
\]
То есть значение функции в промежуточной точке не выше соответствующей точки хорды. Если знак поменять на противоположный, получим вогнутость. Здесь сравниваются конкретные численные значения и их среднее.

\ask{Следует ли из выпуклых предпочтений выпуклая функция полезности? Обязана ли представляющая функция быть выпуклой? Можно ли всегда выбрать выпуклое представление?}

Начнём с $u(x)=x_1+x_2$. Она аффинна: в определении выполняется равенство, поэтому функция одновременно выпукла и вогнута. Её предпочтения выпуклы: смеси безразличных наборов сохраняют сумму.

Теперь сделаем строго возрастающее преобразование $v(x)=e^{x_1+x_2}$. Предпочтения не изменились, но численная кривизна изменилась: новая функция выпукла и уже не вогнута. Таким образом, вогнутость конкретной записи не является порядковым свойством.

Важно заметить предел этого примера: обе функции выпуклы. Поэтому он не доказывает ни существования выпуклого представления для всех выпуклых предпочтений, ни обязательной выпуклости всякого представления. Например, ещё одно представление того же порядка, $w(x)=-e^{-(x_1+x_2)}$, вогнуто и не выпукло. Все три функции одинаково ранжируют наборы.

Сейчас дадим более сильный контрпример к существованию выпуклого представления и затем выделим правильное, порядковое условие.

\slide{42}{Convex preferences and quasiconcavity}
Возьмём предпочтения $u(x)=\sqrt{x_1x_2}$. Они выпуклы. Наборы $a=(1,4)$ и $b=(4,1)$ безразличны: полезность каждого равна двум. Но их среднее $m=(2.5,2.5)$ строго лучше: полезность $2.5$.

Предположим, существует какая-то другая выпуклая функция $g$, представляющая эти же предпочтения. Из безразличия необходимо $g(a)=g(b)$. Выпуклость функции дала бы
\[
g(m)\le\tfrac12g(a)+\tfrac12g(b)=g(a).
\]
Но строгое предпочтение среднего требует $g(m)>g(a)$. Противоречие. Значит, выпуклые предпочтения могут вообще не иметь выпуклого представления полезностью.

Какое же условие правильное? Квазивогнутость:
\[
u(\lambda x+(1-\lambda)y)\ge\min\{u(x),u(y)\}.
\]
Мы сравниваем смесь не со средним двух численных оценок, а с худшей из них. Это условие в точности говорит, что два набора, каждый не хуже общего уровня, дают смесь не хуже того же уровня. Поэтому выпуклость предпочтений эквивалентна квазивогнутости представляющей полезности.

Квазивогнутость сохраняется при строго возрастающих преобразованиях. Вогнутость функции достаточна для неё, но сильнее необходимого. Выпуклость функции сама по себе не является правильным аналогом выпуклости предпочтений: например, $x_1^2+x_2^2$ на положительном ортанте выпукла, но среднее $(1,3)$ и $(3,1)$ имеет значение 8 вместо 10 у каждого конца, то есть предпочтения не выпуклы.

Итак, различаем три объекта: выпуклая функция, выпуклые предпочтения и квазивогнутая полезность. Последние два связаны эквивалентностью. Линейная функция и экспонента показывают, почему числовую кривизну нельзя читать как неизменное свойство ранжирования.

\slide{43}{Perfect substitutes: a straight-line ranking}
Первый стандартный пример --- совершенные заменители:
\[
u(x_1,x_2)=x_1+x_2.
\]
Одна единица первого блага может заменить одну единицу второго без изменения ранга набора. Кривые безразличия --- прямые с наклоном минус один.

\pictureof{43}

Наборы $(4,0)$, $(0,4)$ и $(2,2)$ имеют одинаковую полезность. Значит, предпочтения выпуклы, но не строго выпуклы: смешение двух разных безразличных наборов не даёт строгого улучшения. При этом они сильно монотонны: если прибавить хотя бы одно благо без потерь других, сумма строго увеличится.

Экономическая интерпретация --- два полностью взаимозаменяемых источника одной и той же услуги. Это модель, а не утверждение, что любые два бренда буквально одинаковы. Если единицы благ дают разную эффективность, можно поставить разные положительные коэффициенты и получить другой постоянный обменный коэффициент.

\slide{44}{Perfect complements: useful in fixed proportions}
Второй пример --- совершенные комплементы:
\[
u(x_1,x_2)=\min\{x_1,x_2\}.
\]
Блага полезны в комплектах один к одному. Лишние единицы того блага, которого уже слишком много относительно второго, не улучшают положение. Кривые безразличия имеют излом.

\pictureof{44}

$(2,2)$ и $(5,2)$ безразличны: в обоих случаях можно использовать только два полных комплекта. Но если немного увеличить оба количества, полезность вырастет. Поэтому есть монотонность в нашем смысле и локальная ненасыщаемость, но нет сильной монотонности.

Выпуклость есть. Если оба набора обеспечивают хотя бы $t$ комплектов, каждая их координата не меньше $t$; смесь сохраняет это свойство. Строгой выпуклости нет: среднее $(2,2)$ и $(5,2)$ по-прежнему даёт только два комплекта.

Позднее этот пример напомнит, почему нельзя всегда искать гладкое касание. В изломе единственного гладкого наклона кривой нет, но задача выбора совершенно осмысленна.

\slide{45}{Cobb--Douglas: mixing can help}
Третий пример --- $u(x)=\sqrt{x_1x_2}$. У наборов $(1,4)$ и $(4,1)$ полезность равна двум. У их среднего $(2.5,2.5)$ она равна $2.5$. Здесь смешение двух равноценных крайних сочетаний строго помогает.

\pictureof{45}

На строго положительном ортанте эти предпочтения сильно монотонны и строго выпуклы. При фиксированном положительном уровне кривые имеют вид $x_2=k/x_1$, где $k$ --- квадрат значения полезности. Это привычная гладкая картинка, но сейчас мы видим, какие именно экономические свойства за ней стоят.

С областью определения нужно быть аккуратным. На всём $\R_+^2$ есть оси: если одна координата нулевая, полезность ноль. Наборы $(1,0)$ и $(2,0)$ и их среднее равноценны. Поэтому строгие свойства, сформулированные на положительном ортанте, нельзя без оговорки перенести на границу.

Сравним три примера: прямые заменителей допускают безразличие к смешению; комплементы допускают бесполезный избыток одного блага; Кобб--Дуглас на положительном пространстве даёт и строгую монотонность, и строгий выигрыш от смешения равноценных разных наборов. Функциональная форма всегда несёт содержательные ограничения.

\section{Сепарабельность и квазилинейность: слайды 46--49}
\slide{46}{Separability: a ranking inside a larger bundle}
\reserve
В прикладной модели мы часто хотим изучать небольшую группу товаров, не описывая во всех деталях остальное потребление. Разделим набор на $y$ и $z$. Пусть $y$ --- еда, а $z$ --- жильё и транспорт.

Скажем, что предпочтения над $Y$ не зависят от $Z$, если ранжирование двух пищевых наборов одинаково при любом общем дополнении $z$:
\[
(y,z)\succeq(y',z)
\quad\Longleftrightarrow\quad
(y,z')\succeq(y',z').
\]
Внутри каждого сравнения дополнение одинаковое: мы не сопоставляем один пищевой набор в дорогом доме с другим в дешёвом. Мы спрашиваем, изменится ли рейтинг еды, если поменять общее жильё у обоих вариантов.

Это не означает, что покупка дорогой машины не повлияет на расходы на еду. Бюджет может измениться, хотя рейтинг двух конкретных пищевых наборов сохранится. Предпочтения и ограничения опять нужно различать.

Условие также несимметрично. Можно всегда предпочитать стейк курице при любом вине, но выбирать красное вино к стейку и белое к курице. Например, значения $4,3$ для стейка с красным и белым и $1,2$ для курицы с красным и белым дают именно такую картину. Рейтинг еды независим от вина, рейтинг вина зависит от еды.

\slide{47}{Separability: summarise a group of goods}
\reserve
Смысл результата очень практический: всю внутреннюю сложность пищевого набора можно передать общей функции полезности через один индекс $v(y)$. Сначала считаем пищевой рейтинг, затем соединяем его с остальным потреблением:
\[
u(y,z)=U(v(y),z).
\]

Для каждого фиксированного $z$ агрегатор $U$ строго возрастает по первому аргументу. Иначе более высокий пищевой рейтинг мог бы вдруг превратиться в более низкий общий рейтинг при данном жилье, а это противоречит предпосылке независимости.

В примере на слайде
\[
v(y)=\sqrt{y_1y_2},\qquad U(v,z)=v+\log(1+z).
\]
У пищевого набора $(1,4)$ индекс два, у $(1,1)$ --- один. Какое бы одинаковое $z$ мы ни добавили, первый останется лучше. Это и есть содержательная идея: подробности еды входят в общую оценку через один ранжирующий показатель.

Аддитивная запись здесь лишь пример. Общая теорема не требует, чтобы $U$ была суммой пищевой и непищевой частей. Индекс также не становится объективной единицей счастья или обязательной денежной оценкой. Это порядковое представление внутри группы.

Формальное построение устроено так: фиксируем одно $z_0$, называем $v(y)=u(y,z_0)$ и используем независимость, чтобы перенести равенства и порядок этих индексов на остальные $z$. Полное доказательство оставим приложению. Для прикладной интуиции достаточно понять, какая информация о $y$ нужна агрегатору и что именно предпосылка исключает.

\slide{48}{Quasilinearity: a numeraire plus a valuation}
\reserve
Ещё одна часто используемая форма --- квазилинейность:
\[
u(a,y)=a+v(y),\qquad a\ge0.
\]
$a$ --- нумерар: деньги, оставшиеся на другие расходы, или другое благо, в единицах которого удобно выражать оценки. $v(y)$ описывает ценность группы $y$ на выбранной шкале.

Важное свойство: если к обоим сравниваемым вариантам добавить одну и ту же допустимую сумму нумерара, порядок не меняется:
\[
(a,y)\succeq(a',y')\iff(a+t,y)\succeq(a'+t,y').
\]
Поэтому компенсация за изменение $y$ не зависит от начального уровня $a$ в смысле этой структуры предпочтений. Это более сильное ограничение, чем просто независимость пищевого ранжирования от жилья.

Такая модель удобна, когда хочется сфокусироваться на небольшой части решения. Но удобство оплачивается предпосылкой о роли денег. В частности, привычный вывод об отсутствии эффекта дохода для $y$ имеет оговорку: ограничение на неотрицательность нумерара не должно мешать желаемому выбору. Сейчас увидим это в числах.

\slide{49}{No income effects: mind the boundary}
\reserve
Пусть $u(a,y)=a+4\sqrt y$, оба блага стоят единицу, а бюджет равен $m$. Поскольку увеличение $a$ всегда полезно, при фиксированном $y$ остаток бюджета идёт в $a$: $a=m-y$. Получаем
\[
\max_{0\le y\le m}\bigl(m-y+4\sqrt y\bigr).
\]
Без верхней границы $y\le m$ лучший уровень равен четырём. Это видно даже без производных:
\[
-y+4\sqrt y=4-(\sqrt y-2)^2.
\]
Максимум достигается при $\sqrt y=2$.

При доходе восемь выбираем четыре единицы $y$, остаётся четыре единицы $a$. При доходе двенадцать $y$ по-прежнему четыре, а $a$ уже восемь. Дополнительный доход ушёл в нумерар. Но при доходе два купить четыре единицы $y$ невозможно. Выбор будет $y=2$, $a=0$.

Итак, $y^*=\min\{m,4\}$. При $m>4$ спрос на $y$ не зависит от дохода. При $m<4$ ограничение связывает выбор, и доход имеет значение. Поэтому фраза «у квазилинейной полезности нет эффектов дохода» без оговорки о границе может ввести в заблуждение.

Мы использовали конкретную простую задачу как иллюстрацию. Систематические методы решения общей задачи потребителя будут на следующей лекции.

\section{Бюджет и общая картина выбора: слайды 50--54}
\slide{50}{What can they do? The budget set}
До сих пор мы в основном обсуждали, чего люди хотят. Теперь добавим возможности. Если цены положительны, доход равен $m$, а количества неотрицательны, бюджетное множество двух благ имеет вид
\[
B(p,m)=\{x\in\R_+^2:p_1x_1+p_2x_2\le m\}.
\]
Это всё, что можно купить, а не только те наборы, которые тратят весь доход. Бюджетная прямая --- граница $p_1x_1+p_2x_2=m$. Выразим второе благо:
\[
x_2=\frac{m}{p_2}-\frac{p_1}{p_2}x_1.
\]
Пересечения с осями --- $m/p_1$ и $m/p_2$, наклон --- $-p_1/p_2$. Если вы покупаете дополнительную единицу первого блага при фиксированном бюджете, приходится отказаться от $p_1/p_2$ единицы второго. Это рыночный обменный коэффициент.

Не смешивайте эту линию с кривой безразличия. Одна описывает цены и ресурсы, другая --- желания. У разных людей с одинаковым доходом и ценами может быть одинаковый бюджет и совершенно разные предпочтения.

\ask{Пусть $p_1=2$, $p_2=1$, $m=12$. Где пересечения с осями?}
\answer{Шесть единиц первого блага или двенадцать второго. Доступен весь треугольник под прямой в неотрицательном квадранте. Наклон равен минус двум.}

\slide{51}{Budget geometry: income versus a price change}
Посмотрим на два разных изменения возможностей. Если доход вырос при неизменных ценах, оба пересечения с осями увеличиваются пропорционально, а наклон не меняется. Бюджетная прямая сдвигается параллельно наружу. Старые доступные наборы остаются доступными, добавляются новые.

Если выросла только цена первого блага при неизменных втором ценнике и доходе, вертикальное пересечение $m/p_2$ сохраняется, а горизонтальное $m/p_1$ уменьшается. Линия поворачивается внутрь и становится круче. Это не параллельный сдвиг.

\pictureof{51}

\ask{Вернёмся к примеру $p_1=2$, $p_2=1$, $m=12$. Что произойдёт, если доход вырастет до 18? А если вместо этого первая цена вырастет до трёх при прежнем доходе?}
\answer{В первом случае пересечения станут $9$ и $18$, наклон остаётся минус два. Во втором --- $4$ и $12$, наклон минус три.}

Пока мы сделали выводы только о множестве возможностей. Из расширения бюджета не следует, что потребление каждого блага обязательно вырастет. Для такого прогноза нужно соединить бюджет с предпочтениями и решить задачу. Ресурсы и вкусы --- разные части модели.

\slide{52}{Wants and possibilities in one picture}
Теперь объединим две половины. На рисунке полезность $u(x)=\sqrt{x_1x_2}$, цены обоих благ равны единице, доход восемь. Значит, доступны наборы $x_1+x_2\le8$.

\pictureof{52}

Сначала виден бюджет. Затем появляются кривые безразличия с уровнями два, четыре и шесть. Уровень шесть хотелось бы достичь, но он целиком вне доступного множества. Уровень два достижим, однако есть доступные наборы лучше. Лучшая доступная кривая касается бюджета в $(4,4)$, где полезность равна четырём.

При желании это можно проверить без условий первого порядка. Для любого доступного набора
\[
x_1x_2\le\left(\frac{x_1+x_2}{2}\right)^2\le16,
\]
поэтому $\sqrt{x_1x_2}\le4$. Равенство требует $x_1=x_2=4$. То есть точка на картинке действительно глобально лучшая, а не просто похожа на неё.

Но не превращайте это в правило «всегда ищем касание». У совершенных комплементов есть изломы, у заменителей могут возникать угловые решения, а свойства существования и единственности требуют отдельного разговора. Сегодня главное --- увидеть, как желания и возможности вместе определяют выбор. Следующая лекция систематизирует решение.

\slide{53}{Exit ticket}
Проверим, какие различия остались в голове. Почему возрастающее преобразование может изменить вогнутость, но не предпочтения? Потому что порядок сохраняется, а численные разности и кривизна --- нет. Достаточно сравнить $x_1+x_2$ и $e^{x_1+x_2}$.

Пример рациональных, но не сильно монотонных предпочтений --- $u=\min\{x_1,x_2\}$. Полнота и транзитивность следуют из числового представления, но дополнительное количество избыточного блага может ничего не дать.

Обещает ли выпуклость строгое улучшение от любой смеси? Нет: она сохраняет слабое сравнение с общим исходным уровнем. Сепарабельность исключает эффекты дохода? Тоже нет: она касается ранжирования внутри группы. Например, $\log y+\log z$ сепарабельна, но при единичных ценах даёт $y^*=m/2$.

Наконец, в чём граничная оговорка квазилинейности? Если неотрицательность нумерара связывает решение, спрос на другую группу может зависеть от дохода. В нашем примере это $y^*=\min\{m,4\}$.

\cue{Не читать готовые ответы подряд до ответов группы. Эти абзацы --- разбор для преподавателя. Если резервные темы были пропущены, последние два вопроса можно обсудить как анонс или не требовать ответа на ещё не пройденное.}

\slide{54}{What to keep, and what to read}
Сегодня мы отделили три вещи: альтернативы, их ранжирование и доступное меню. Определили рациональность как полноту и транзитивность, связали предпочтения с правилом выбора и разобрались с полезностью как числовым представлением. Затем посмотрели, что означают непрерывность, монотонность, LNS и выпуклость, и собрали бюджет с желаниями на одном рисунке.

Основное чтение --- Hal Varian, \emph{Intermediate Microeconomics: A Modern Approach}, девятое издание 2014 года: глава 2 \emph{Budget Constraint}, глава 3 \emph{Preferences}, глава 4 \emph{Utility}. На следующую встречу важно уверенно читать определения и три стандартные геометрии, а не только узнавать знакомую формулу.

Дополнительное продвинутое чтение --- Mas-Colell, Whinston and Green, \emph{Microeconomic Theory} 1995 года: глава 1 \emph{Preference and Choice}, а также глава 3 \emph{Classical Demand Theory}, разделы 3.B--3.C о свойствах предпочтений и полезности. Целиком всю дуальность из этой главы сейчас читать не требуется.

И начните замечать загадки для проекта. Хорошее наблюдение может быть совсем бытовым. Главное --- увидеть выбирающего участника, его ограничения и механизм, который можно поставить под проверку. Следующая лекция превращает нашу общую картинку в задачу максимизации и функции спроса.

\newpage
\section{Полные доказательства и дополнительные разборы: слайды 55--64}
Эта часть нужна прежде всего для подготовки и ответов на вопросы. В аудитории можно использовать отдельные аргументы по ссылкам из презентации. Она не является дополнительным обязательным списком доказательств, которые нужно успеть произнести за две пары.

\slide{55}{Why rational preferences imply WARP}
Пусть $\succeq$ полно и транзитивно, а $C(A)$ выбирает все лучшие элементы меню. Возьмём $x,y\in A\cap B$, причём $x\in C(A)$ и $y\in C(B)$. Нужно показать два обратных включения в выбранные множества.

Поскольку $y$ доступен в $A$, а $x$ там лучший, имеем $x\succeq y$. Поскольку $x$ доступен в $B$, а $y$ там лучший, имеем $y\succeq x$. Значит, $x\sim y$.

Теперь возьмём произвольный $z\in B$. Оптимальность $y$ даёт $y\succeq z$. Вместе с $x\succeq y$ транзитивность даёт $x\succeq z$. Это верно для всех $z\in B$, поэтому $x\in C(B)$. Симметрично для любого $w\in A$ из $y\succeq x\succeq w$ получаем $y\succeq w$, значит, $y\in C(A)$.

Аксиома доказана. В самом этом переходе конечность не использовалась: если лучшие элементы существуют, аргумент работает и на бесконечном меню. Конечность нужна в формулировке общей теоремы, чтобы не скрыть проблему существования выбора. В выбранных точках мы постоянно использовали транзитивность, а не вычислительную способность агента.

\cue{Вернуться к \hyperlink{script-25}{слайду 25}, если доказательство показано как отступление.}

\slide{56}{Reconstructing the relation from choice}
Теперь обратное направление. Пусть $X$ конечно, $C$ определено и непусто на всех непустых меню и выполняет WARP. Определим
\[
x\succeq y\quad\Longleftrightarrow\quad x\in C(\{x,y\}).
\]
Нужно доказать полноту, транзитивность и совпадение $C$ с выбором лучших элементов этого отношения.

\textbf{Полнота.} В парном меню что-то выбрано. Если выбран $x$, получаем $x\succeq y$; если $y$, то $y\succeq x$. Если оба, получаем безразличие. Для одноэлементного меню выбран его единственный элемент, так что есть и рефлексивность.

\textbf{Транзитивность.} Пусть $x\succeq y$ и $y\succeq z$. Если какие-то элементы совпадают, нужный вывод непосредственный. Для трёх разных альтернатив рассмотрим $T=\{x,y,z\}$ и выберем $w\in C(T)$.

Если $w=x$, уже знаем, что $x$ выбран в тройке. Если $w=y$, сравним тройку с парой $\{x,y\}$, из которой $x$ выбран по предположению. WARP переносит $x$ в $C(T)$. Если $w=z$, сначала сравним $T$ с $\{y,z\}$: поскольку $y$ выбран в паре, WARP переносит $y$ в тройку. Затем предыдущий шаг переносит туда $x$.

Итак, во всех случаях $x\in C(T)$. В паре $\{x,z\}$ есть выбранный элемент. Если это $x$, получили $x\succeq z$. Если выбран $z$, WARP между этой парой и $T$ заставляет включить $x$ в выбранное множество пары. Снова $x\succeq z$. Транзитивность доказана.

\textbf{Совпадение правил.} Сначала пусть $x\in C(A)$. Для каждого $y\in A$ рассмотрим пару $\{x,y\}$. Если в ней выбран $x$, нужное сравнение уже есть. Если выбран $y$, WARP с меню $A$ заставляет выбрать в паре и $x$. Значит, $x\succeq y$ для всех $y\in A$: выбранный элемент действительно лучший.

Наоборот, пусть $x\in A$ и $x\succeq y$ для всех $y\in A$. Возьмём какой-нибудь $y\in C(A)$, что возможно по непустоте. По построению $x$ выбран из $\{x,y\}$. WARP между парой и меню $A$ даёт $x\in C(A)$. Следовательно, выбраны ровно все лучшие элементы.

Обратите внимание, сколько информации использовано. Нам нужны парные меню, тройки, непустота и само исходное меню $A$. Именно поэтому переносить результат на произвольный обрывок эмпирических данных нельзя автоматически.

\slide{57}{Limited data can hide a cycle}
Пусть наблюдаем только три парных меню:
\[
C(\{s,p\})=\{s\},\quad
C(\{p,r\})=\{p\},\quad
C(\{s,r\})=\{r\}.
\]
Это соответствует строгим сравнениям $s\succ p$, $p\succ r$, $r\succ s$. Единого полного транзитивного отношения, порождающего именно эти полные выбранные множества, нет.

Однако разные наблюдаемые меню пересекаются лишь в одном элементе. Между ними нет двух различных общих альтернатив, на которых WARP мог бы выявить обратный выбор. При сравнении меню с самим собой выбранное множество тоже не даёт нарушения. Поэтому на этом ограниченном домене тест проходит, а глобальная рационализируемость отсутствует.

Это хороший пример разницы между необходимым условием на выборке и достаточным условием для полного объекта. Если добавить тройное меню, согласованно продолжить правило при сохранении WARP не получится.

И опять оговорка о данных: мы предполагаем, что выбранные множества действительно одноэлементны. Если наблюдали лишь по одному реализованному действию, универсальное безразличие с разным разрешением ничьей могло бы объяснить записи. Строгий цикл тогда нельзя вывести без дополнительной информации.

\slide{58}{Why lexicographic utility does not exist}
Предположим для противоречия, что лексикографические предпочтения на $[0,1]^2$ представляет вещественная функция $u$.

Для каждого $t\in[0,1]$ имеем $(t,1)\succ(t,0)$, поэтому
\[
u(t,1)>u(t,0).
\]
Значит, открытый интервал $I_t=(u(t,0),u(t,1))$ непуст. Если $s>t$, то по лексикографическому правилу $(s,0)\succ(t,1)$, следовательно, $u(s,0)>u(t,1)$. Таким образом, интервалы $I_s$ и $I_t$ не пересекаются.

В каждом непустом открытом вещественном интервале есть рациональное число. Выберем по одному такому числу из каждого $I_t$. Можно делать это совершенно определённо: заранее перечислить рациональные и брать первое из попавших в интервал. Разные непересекающиеся интервалы дадут разные рациональные числа.

Получили инъекцию из несчётного множества $[0,1]$ в счётное множество рациональных чисел. Это невозможно. Значит, никакого вещественного представления не существует.

Доказательство сильнее утверждения об отсутствии непрерывной полезности. Мы исключили любую вещественную функцию, даже разрывную. А рациональность предпочтений при этом не оспаривалась. Пример показывает реальное ограничение числового языка на бесконечных пространствах.

\slide{59}{Construction under monotonicity}
Пусть на $\R_+^L$ предпочтения полны, транзитивны, непрерывны и монотонны в смысле строгого улучшения при увеличении всех координат. Построим непрерывную полезность через диагональный эквивалент.

Сначала заметим: если $x\ge y$, то $x+\varepsilon e\gg y$ при любом $\varepsilon>0$, где $e=(1,\ldots,1)$. Монотонность даёт $x+\varepsilon e\succ y$, а непрерывность при $\varepsilon\downarrow0$ даёт $x\succeq y$. Для фиксированного $x$ возьмём $M=\max_i x_i$. Тогда $0\le x\le Me$, откуда
\[
x\succeq0,\qquad Me\succeq x.
\]

На отрезке $[0,M]$ определим два множества коэффициентов:
\[
A_- =\{a:x\succeq ae\},\qquad
A_+ =\{a:ae\succeq x\}.
\]
Первое содержит ноль, второе --- $M$. По полноте их объединение --- весь отрезок. По непрерывности оба замкнуты: предел коэффициентов, сохраняющих соответствующее слабое сравнение, тоже сохраняет его. Связный отрезок нельзя разбить на два непустых непересекающихся замкнутых множества. Значит, их пересечение непусто: существует $a$ с $x\sim ae$. При $x=0$ вывод непосредственный.

Эквивалент единственен. Если $a>b$, то $ae\gg be$, поэтому $ae\succ be$. Эти два диагональных набора не могут оба быть безразличны одному $x$, иначе транзитивность дала бы их безразличие друг другу.

Определим $u(x)=a$, где $x\sim ae$. Если $u(x)\ge u(y)$, соответствующий диагональный набор $x$ не хуже диагонального набора $y$, и транзитивность переносит сравнение на исходные наборы. Если $u(x)<u(y)$, получаем строгое обратное сравнение. Поэтому
\[
x\succeq y\iff u(x)\ge u(y).
\]

Осталось проверить непрерывность, а не только представление. Пусть $x^n\to x$ и $a_n=u(x^n)$. Последовательность наборов ограничена по координатам, поэтому существует общий $M$ с $0\le a_n\le M$. Возьмём любую сходящуюся подпоследовательность коэффициентов, $a_{n_k}\to a$. Для каждого её члена $x^{n_k}\sim a_{n_k}e$. Непрерывность предпочтений в обе стороны даёт $x\sim ae$. По единственности $a=u(x)$.

Таким образом, всякая сходящаяся подпоследовательность коэффициентов имеет один и тот же предел. Поскольку коэффициенты лежат в компактном отрезке, это означает $a_n\to u(x)$. Непрерывность построенной функции доказана.

Мы доказали полезный частный случай общей теоремы. Дополнительная монотонность позволила выбрать диагональ и однозначный коэффициент. Без неё общее существование непрерывной полезности на данном пространстве требует другого аргумента.

\slide{60}{Separability in one construction}
Пусть $u(y,z)$ представляет предпочтения, а ранжирование $Y$ не зависит от $Z$. Фиксируем $z_0$ и полагаем $v(y)=u(y,z_0)$.

Для каждого достигаемого значения $t\in v(Y)$ выберем представителя $\hat y(t)$ с $v(\hat y(t))=t$. Определим
\[
U(t,z)=u(\hat y(t),z).
\]
Нужно проверить, что выбор представителя не влияет на результат. Если $v(y)=v(y')$, наборы $(y,z_0)$ и $(y',z_0)$ безразличны. Независимость переносит оба слабых сравнения на любое общее $z$. Значит, $u(y,z)=u(y',z)$. Следовательно, $U$ определена корректно.

Теперь $y$ и $\hat y(v(y))$ имеют один пищевой индекс, а значит, безразличны при любом общем $z$. Поэтому
\[
U(v(y),z)=u(\hat y(v(y)),z)=u(y,z).
\]
Если $t>t'$, при $z_0$ соответствующие представители строго ранжированы. Независимость сохраняет строгое сравнение при любом $z$, откуда $U(t,z)>U(t',z)$. Так получаем строгое возрастание агрегатора по первому аргументу.

Обратное направление проще: если $u(y,z)=U(v(y),z)$ и $U$ строго возрастает по $v$, то при любом фиксированном $z$ ранжирование $y,y'$ совпадает с ранжированием $v(y),v(y')$. Последнее не зависит от $z$, значит, независимость выполнена.

Область построенного $U$ --- $v(Y)\times Z$, а не автоматически вся $\R\times Z$. Продолжать агрегатор на недостигаемые значения индекса для этого результата не требуется.

\slide{61}{Primitives behind quasilinearity}
Покажем достаточный набор исходных условий для записи $u(a,y)=a+v(y)$. Пусть предпочтения рациональны на $\R_+\times Y$, существует худший элемент $\bar y$, а нумерар при этом элементе правильно ранжируется по количеству:
\[
(a,\bar y)\succeq(a',\bar y)\iff a\ge a'.
\]
Пусть для каждого $y$ существует неотрицательная компенсация $t$ с $(0,y)\sim(t,\bar y)$. И пусть добавление одинакового допустимого количества нумерара к обеим альтернативам сохраняет сравнение в обе стороны.

Определим $v(y)$ как компенсирующее количество: $(0,y)\sim(v(y),\bar y)$. Оно единственно, потому что разные количества нумерара при $\bar y$ строго ранжируются. Добавив $a$, получаем
\[
(a,y)\sim(a+v(y),\bar y).
\]
Аналогично $(a',y')\sim(a'+v(y'),\bar y)$. По транзитивности сравнить исходные наборы --- то же, что сравнить эти два количества нумерара при худшем $\bar y$. Значит,
\[
(a,y)\succeq(a',y')\iff a+v(y)\ge a'+v(y').
\]
Представление доказано.

Важно сохранять условия об области и компенсации. Из произвольной квазилинейной формулы на любом множестве $Y$ не следует автоматически существование худшего элемента $\bar y$ и неотрицательной компенсации для каждого $y$. Здесь они являются предпосылками построения, а не бесплатными выводами из названия формы.

\slide{62}{LNS and budget exhaustion}
Пусть $p\gg0$, предпочтения локально ненасыщаемы, а $x^*$ --- оптимальное решение при бюджете $p\cdot x\le m$. Предположим, что оно оставляет положительный запас денег:
\[
\delta=m-p\cdot x^*>0.
\]
Для достаточно близкого $y$ бюджет ещё выполняется. Точнее, по неравенству Коши--Буняковского при $\|y-x^*\|<\delta/\|p\|$ получаем
\[
p\cdot y\le p\cdot x^*+\|p\|\,\|y-x^*\|<m.
\]
LNS гарантирует в этой окрестности допустимый для пространства потребления набор $y\succ x^*$. Он доступен и строго лучше, что противоречит оптимальности $x^*$.

Значит, существующее оптимальное решение исчерпывает бюджет: $p\cdot x^*=m$. В доказательстве нигде не устанавливалось само существование $x^*$. Также не требовалось, чтобы любая прибавка любого отдельного блага улучшала положение: достаточно некоторого локального улучшения.

Экономическая интерпретация относится к правильно заданному набору благ. Если можно сберегать, сбережения тоже должны входить в модель как вариант использования ресурсов. Бюджетное равенство не является утверждением, что реальные люди обязаны тратить всё на текущее потребление.

\slide{63}{Increasing transformations and risk}
Возьмём результат $A=2$ наверняка и случайный результат $B$, равный нулю или трём с вероятностями по половине. По ожидаемому исходному числу
\[
\mathbb E[A]=2>1.5=\mathbb E[B].
\]
Теперь применим строго возрастающую экспоненту к каждому исходу до усреднения:
\[
\mathbb E[e^A]=e^2\approx7.39,
\qquad
\mathbb E[e^B]=\frac{1+e^3}{2}\approx10.54.
\]
Ранжирование изменилось.

Это не противоречит порядковой инвариантности для детерминированного выбора. Мы не просто применили возрастающую функцию к уже готовому показателю $\mathbb E[A]$ или $\mathbb E[B]$. Мы поменяли численную оценку исходов до операции ожидания. В общем случае $\mathbb E[f(X)]$ не равно $f(\mathbb E[X])$.

Поэтому полезность отдельных результатов и ожидаемая полезность лотерей требуют разного обращения со шкалой. В блоке о риске появятся дополнительные аксиомы и более узкий класс допустимых преобразований ожидаемой полезности. Похожая осторожность нужна в RL: нелинейно преобразовать каждую награду перед суммированием --- не то же самое, что монотонно преобразовать всю итоговую целевую функцию.

\slide{64}{Answers to the relation exercise}
Короткая памятка для возвращения к упражнению. Сумма координат задаёт полное транзитивное отношение. Правило большинства трёх координат полно, но допускает цикл на $(3,2,1)$, $(2,1,3)$ и $(1,3,2)$. Покомпонентный порядок транзитивен, но неполон на парах с разнонаправленными преимуществами.

Это три разных структуры. В первом случае можно сразу использовать вещественную сумму как полезность. Во втором вещественное представление всех сравнений невозможно уже из-за нетранзитивности. В третьем --- из-за неполноты: вещественные числа всегда сравнимы, а исходное отношение оставляет некоторые пары несравнимыми.

Парадокс Кондорсе на слайде 20 даёт интерпретацию второго случая через агрегирование трёх рациональных людей. Он показывает, почему переход от индивидуальных сравнений к коллективному результату требует отдельного анализа.

\section{Последняя проверка перед занятием}
\begin{itemize}[leftmargin=*]
\item Основной сюжет: чего хотят, что могут, как соединяются предпочтения и бюджет.
\item Не путать безразличие с несравнимостью, реализованную покупку со всем выбранным множеством, представление с существованием максимума.
\item Отдельно проговорить смысл слова monotone: в нашей терминологии оно изначально относится к строгому росту всех координат. Именно поэтому нужен небольшой аргумент с непрерывностью.
\item На ресторанном примере сначала собрать объяснения, затем открыть условия идентичности общих альтернатив и единственного фаворита.
\item На Кондорсе каждый избиратель транзитивен, цикл появляется только у большинства. Это не общая теорема о неизбежности циклов.
\item В судебном кейсе говорить об обвинениях из жалобы 2023 года; учебные числа не представлять как реальные комиссии.
\item Домашние выдаются после блоков 3, 5, 7. Проект сохраняет свои вехи после 2, 4 и 6 и презентацию на семинаре 7.
\item При нехватке времени сокращать резервные темы, сохраняя бюджет и общую картинку в конце.
\end{itemize}

\section{Источники для подготовки}
Основная последовательность математического материала адаптирована из предоставленных преподавателем документов \emph{Lecture 8: Consumers!}, печатные страницы 119--132, и \emph{Lecture 9: Properties of Preferences}, страницы 133--148. Авторство в заголовках и метаданных этих файлов не указано; непроверенная атрибуция не добавляется. Текст выступления, связки, числовые примеры и диаграммы подготовлены для текущего курса.

Назначенное студентам чтение: Hal R. Varian, \emph{Intermediate Microeconomics: A Modern Approach}, 9th edition, 2014, главы 2--4. Дополнительно: Andreu Mas-Colell, Michael D. Whinston, Jerry R. Green, \emph{Microeconomic Theory}, 1995, глава 1 и разделы 3.B--3.C.

Для кейса приведена \href{https://www.ftc.gov/news-events/news/press-releases/2023/09/ftc-sues-amazon-illegally-maintaining-monopoly-power}{официальная публикация FTC о жалобе против Amazon 2023 года}. Текущий статус дела в рамках этой подготовки не проверен. Информация о недоступности онлайн-проверки относится к подготовке, а не является частью речи перед аудиторией.

\end{document}
